\documentclass[10pt,a4paper]{amsart}
\usepackage[margin=19mm]{geometry}
\usepackage{amsmath,amssymb,mathtools}
\usepackage[scr=rsfs,cal=euler]{mathalfa}
\usepackage{xfrac}
\usepackage[unicode,hidelinks,bookmarks=false]{hyperref}
\newcommand{\ie}{i.e.\ }
\newcommand{\cf}{cf.\ }
\newcommand{\resp}{respectively\ }
\newcommand{\Subsection}{\subsection}
\providecommand{\nobreakdash}{\nobreak\hbox{-}}
\setlength{\emergencystretch}{2em}
\multlinegap0pt
\usepackage{bm}
\usepackage{mathtools}
\newtheorem{lemma}{Lemma}
\newcommand{\aX}{\mathbb{A}[X]}
\title{Computing Certificates in Archimedean Univariate Saturated Quadratic Modules}
\author{Jose Abel Castellanos-Joo, Deepak Kapur}
\date{Source: arXiv:2605.18980v1 (CC BY 4.0)}
\begin{document}
\maketitle

\noindent\emph{Source and licence.} Computing Certificates in Archimedean Univariate Saturated Quadratic Modules. Version arXiv:2605.18980v1 (CC BY 4.0), distributed by arXiv under the Creative Commons Attribution 4.0 International licence, http://creativecommons.org/licenses/by/4.0/, as recorded in the arXiv licence metadata for this exact version and shown on the abstract page https://arxiv.org/abs/2605.18980v1. Full licence notice: \texttt{licenses/arxiv-2605.18980-CC-BY-4.0.txt}.\par\smallskip

\noindent\textit{Editorial scope.} Counted unit: the article's lemma sec:splitting-step-7 (source0.tex lines 2518--2523). The article's complete original proof of it is delivered with it (source0.tex lines 2525--2619, one \texttt{proof} environment of 95 lines). No mathematics is written, repaired or re-derived: the statement and the proof are the article's own lines, in the article's own notation, and no retained line is substituted or re-flowed.  No further source-local context is needed: every cross-reference of every retained line resolves inside the two delivered spans.  Nothing was omitted from any retained span, so the package carries no omission record.  No retained line contains a citation, so the wrapper carries no bibliography and no bibliography entry.  No retained line contains a figure environment, an image-inclusion command or drawing code, so no image is delivered and the package declares no asset dependency.  Editorial formatting only, in addition: an AMS \texttt{amsart} preamble that restates the article's own declarations and macros used by the retained lines, namely the environments \texttt{lemma} on separate counters, together with the standard packages those lines need.  The retained environments print in this document as Lemma 1 in order of appearance, whereas the article prints the same items with the numbering of its own full document.  The complete native source of the article as distributed in the arXiv e-print for this exact version is bundled unchanged as \texttt{sources/200884/source0.tex}.
\begin{lemma} \th\label{sec:splitting-step-7} Let
  $p := \alpha(X-a)(X-b)^{m} + \beta(X-c)^{m}(X-d)$ with
  $m \in 2\mathbb{N} + 1$, $b, c \in \mathbb{A}$ within $(a, d)$, and
  $\alpha := \frac{1}{(d-a)(d-b)^{m}}, \beta := \frac{1}{(a-c)^{m}(a-d)}$. Then
  $p$ is of the form $q \cdot (X-a)(X-d) + 1$ for some $q \in \sum{\aX}^2$.
\end{lemma}

\begin{proof}
  Notice that $p(a) = p(d) = 1$. Using Euclid's division algorithm, the residue
  of $p$ divided by $(X-a)(X-d)$ is at most of degree 1, i.e.,
  $p = q(X-a)(X-d) + r$ with $\deg(r) \leq 1$. Let $r = \gamma_1 X +
  \gamma_2$. If $\gamma_1 = 0$ we are done as $\gamma_2 = 1$. Otherwise,
  $p(a) = r(a) = 1$, similarly $p(d) = r(d) = 1$. Hence, $r(a) = r(d)$, so
  $\gamma_1 a + \gamma_2 = \gamma_1 d + \gamma_2$, which implies that $a =
  d$. This contradicts the assumption that $a < d$.

  Now, we will show $q \in \sum{\aX}^2$ to conclude the proof. We will
  show that the only real roots of the polynomial $p - 1$ are $a, d$
  with multiplicity 1, which shows that $q$ must be a product of complex
  conjugate roots, which means that $q$ is a sum of squares.

  Suppose $q$ has a real root $e$ different from $a$ and $d$. This implies
 $p(e) = 1$. We will prove this leads to a contradiction by showing that $p$
  evaluates to a value greater than $1$ over $(-\infty, a)$ and $(d, \infty)$
  and smaller than $1$ over $(a, d)$.

  First, write $p$ as $p = \frac{p_1}{p_1(d)} + \frac{p_2}{p_2(a)}$ where
  $p_1 = (X-a)(X-b)^{m}$ and $p_2 = (X-c)^m(X-d)$. For the interval $x \in (d,
  \infty)$, we see that $\frac{p_1(x)}{p_1(d)}$ is strictly greater than $1$ and
  $\frac{p_2(x)}{p_2(a)}$ is non-negative; hence $p(x)$ is strictly greater than
  $1$. A similar argument holds for the interval $(-\infty, a)$.

  Now, let us consider the case where $x \in (a, d)$. If $b = c$ we have
 $p = \frac{p_2(a) + p_1(d)}{p_2(a) p_1(d)}(X-b)^{m}\left(X - r\right)$ where
  $r = \frac{a p_2(a) + d p_1(d)}{p_2(a) + p_1(d)}$. We can see that $p$ has at
  most two real zeros which belong to the interval $(a, d)$ as $a < r <
  d$. Since $p_1(d), p_2(a)$ are positive, the leading coefficient of $p$ is
  positive. The minimum value $p^{*}$ of $p$ is attained within
 $\left(a, d\right)$; this implies that $p$ strictly decreases in the
 interval $(a, p^{*})$ and increases strictly in the interval $(p^{*},
  d)$. As $p(a) = p(d) = 1$, we conclude that $p$ reaches a value less than
  $1$ over $(a, d)$.

  Now, let us consider the case where $b < c$. We can check that
  $p(b) = \frac{b-d}{a-d} (\frac{b-c}{a-c})^{m} < 1$ and
  $p(d) = \frac{a-c}{a-d} (\frac{b-c}{b-d})^{m} < 1$. To finish this case, we
  analyze the intervals $(a, b), (b, c)$ and $(c, d)$. Let us consider the
  interval $(a, b)$. The polynomial $\frac{p_1}{p_1(d)}$ evaluates to negative
  values over $(a, b)$ and the polynomial $\frac{p_2}{p_2(a)}$ is strictly
  decreasing over $(a, b)$ with maximum value $1$; thus $p$ evaluates to a value
  less than $1$ over the interval $(a, b)$; a similar reasoning handles the
  interval $(c, d)$.

  Now, we address the interval $(b, c)$. First, we notice that
  \vspace{-5mm}

  \begin{equation}
    \label{sec:addit-constr-from-1}
    \begin{split}
      \frac{dp}{dX}
      &= \frac{1}{p_1(d)}((X-b)^{m} + m (X-a)(X-b)^{m-1})\\
      &+
        \frac{1}{p_2(a)}((X-c)^{m} + m (X-d)(X-c)^{m-1})\\
      &= \frac{(X - b)^{m-1}}{(d-a)(d-b)^{m}}((m+1)X - (b + m a))\\
      &+
        \frac{(X-c)^{m-1}}{(a-d)(a-c)^{m}}((m+1)X - (c + m d))\\
    \end{split}
  \end{equation}

  We want to show $\frac{dp}{dX}$ is strictly increasing over $(b, c)$.  The
  polynomial $q_1 := \frac{(X - b)^{m-1}}{(d-a)(d-b)^{m}}((m+1)X - (b + m a))$
  has at most two real roots, $b$ and $\frac{b + ma}{m+1}$. The last is
 contained in the interval $[a, b]$, which means that the polynomial $q_1$ increases
 strictly for any point after $b$. Similarly, we can show that the
  polynomial $\frac{(X-c)^{m-1}}{(a-d)(a-c)^{m}}((m+1)X - (c + m d))$ is
  strictly increasing for any point before $c$. Hence, from equation
  \eqref{sec:addit-constr-from-1}, we see that $\frac{dp}{dX}$ is strictly
  increasing over $(b, c)$.

  We have
  $\frac{dp}{dX}(b) = \frac{1}{b-c} \frac{m(b-d) + (b - c)}{a-d}
  (\frac{b-c}{a-c})^{m} < 0$ as $b < c$. Similarly,
  $\frac{dp}{dX}(c) = \frac{1}{c-b} \frac{m(a-c) + (b - c)}{a-d}
  (\frac{c-b}{d-b})^{m} > 0$. By the Intermediate Value theorem, we have that
  $\frac{dp}{dX}$ has a zero $\eta$ in $(b, c)$ and it is unique since
  $\frac{dp}{dX}$ is strictly increasing over $(b, c)$. Furthermore, $\eta$
  corresponds to the global minimum value of $p$ over $\mathbb{R}$. Hence
  $p(x) < 1$ for $x \in (b, c)$, as otherwise there would be a local maximum
  value of $p$ in $(b, c)$, which is not the case.

  Finally, let us consider the case where $c < b$. For the interval $(a, c)$,
  $\frac{p_1}{p_1(d)}$ is less than $1$ and $\frac{p_2}{p_2(a)}$
  is always non-negative over $(a, c)$, hence $p$ is less than $1$ over $(a,
  c)$. A similar argument handles the interval $(b, d)$. For the interval $[c,
  b]$ we notice that both $\frac{p_1}{p_1(d)}$ and $\frac{p_2}{p_2(a)}$ are
  non-positive, so $p$ is less than $1$ over $[b, d]$. Therefore, $p$ is less
  than $1$ over $(a, d)$.

  Therefore, $q$ has no real roots and is a product of irreducible quadratics
  of the form $q = \prod_{j}^{}{((X - a_j)^2 + b_j^2)^{e_j}}$. Since the leading
  coefficient of $p$ is positive, we see that $q$ is a sum of squares.
\end{proof}

\end{document}
